%%%%%%%%%%  scan idx 77  |  folio 64  %%%%%%%%%%
(2.8) résulte de (2.31).
\begin{flushright}cqfd\end{flushright}

\textbf{Corollaire 1~:} \emph{On a $c^i(x)=\varphi(C^i(x))$.}

Ceci résulte d'un théorème d'unicité pour les classes de Chern, qui se démontre comme chez Hirzebruch par passage aux variétés de drapeaux. On doit seulement montrer que si $Y$ est une variété de drapeaux associée à un fibré vectoriel sur $X$, alors $G(\underline K(X))\longrightarrow G(\underline K(Y))$ est injectif, ce qui se vérifie par la détermination explicité de $\underline K(Y)$ avec toutes ses structures à partir de $\underline K(X)$ (je ne donnerai pas les détails ici).

\textbf{Corollaire 2~:} \emph{Supposons que pour tout élément $Z\in A(X)$, il existe un morphisme $f$ de $X$ dans un produit $X'$ de grasmaniennes tel que $Z$ soit dans l'image de $f^*$. Alors les homomorphismes canoniques}
\[
 \varphi:A(X)\longrightarrow G(\underline K(X)),\qquad \psi:G(\underline K(X))\longrightarrow A(X),
\]
\emph{vérifient en dimension $i$ les relations suivantes~:}
\begin{equation}\tag{2.35}
 \psi\varphi=(-1)^{i-1}(i-1)!\,\underline 1,\qquad
 \varphi\psi=(-1)^{i-1}(i-1)!\,\underline 1.
\end{equation}

Comme $\varphi$ est surjective, il suffit évidemment de prouver la première formule, et pour celle-ci, on peut supposer que $X$ est un produit de grasmaniennes. Comme alors $\varphi$ est injective, il suffit de prouver la deuxième formule (2.35), qui en vertu du corollaire 1, s'écrit
\[
 \gamma^i(x-\varepsilon(X))\equiv(-1)^{i-1}(i-1)!\,x\qquad\mathrm{mod}\ \underline K(X)^{i+1}
\]
pour tout $x\in\underline K(X)^i$. Comme $\psi$ est \emph{injective}, il suffit de prouver la relation analogue pour $\Ctil(x)\in A(X)^i$, et c'est alors une conséquence de (1.22 bis).

\textbf{Remarque~:} Le corollaire 2 rend plausible la validité de (2.35) dans tous les cas~; il devrait en exister une démonstration directe. On prouve en
